\newpage
\section{Acta Número 19}
\label{sec:acta19}

\noindent \textbf{Fecha de la sesión:}\par
Jueves, 7 de Mayo de 2026

\vspace*{0.5cm}
\noindent \textbf{Datos de la Sesión:}
\vspace{-0.2cm}
\begin{itemize}[noitemsep]
    \item \textbf{Modalidad:} Virtual - Google Meet.
    \item \textbf{Hora de Inicio:} 18:00
    \item \textbf{Hora de Fin:} 20:00
\end{itemize}

\vspace*{0.5cm}
\noindent \textbf{Orden del Día:}
\vspace{-0.2cm}
\begin{itemize}[noitemsep]
    \item Control de Asistencia.
    \item Recepción de reportes de Aseguramiento de Calidad (QA) y revisión de Test Cases y Bugs (Sprints 1, 2 y 3).
    \item Entrega y revisión de documentación técnica integral (Backend, Frontend, Base de Datos, Validaciones y Control de Versiones).
    \item Definición de fechas definitivas de entrega y despliegue para el Sprint 3.
    \item Firmas y Aprobación de los Presentes.
\end{itemize}

\vspace*{0.5cm}
\noindent \textbf{Socios Fundadores Presentes:}
\vspace{-0.2cm}
\begin{enumerate}[noitemsep]
    \item Pardo Pozo Juan Diego (Jefe de Proyecto).
    \item Arnez Jaimes Mauricio.
    \item Quispe Flores Camila Belen.
    \item Ariñez Bautista Jim Gabriel.
    \item Medina Rodríguez Mateo Jhoustin.
    \item Molina Maldonado Tomas Manuel.
    \item Gutierrez Guzmán Angheline.
\end{enumerate}

\newpage
\subsection{Desarrollo del Acta}

\noindent \textbf{Control de Asistencia del Team}\par
Se procedió a la toma de asistencia a la hora acordada, corroborando la presencia de los 7 socios fundadores de la grupo-empresa Byte Busters S.R.L. Al contar con el quórum y la atención de todos los involucrados, se dio inicio formal a la sesión.

\vspace*{0.5cm}
\noindent \textbf{Revisión de Aseguramiento de Calidad (QA): Reportes, Bugs y Test Cases}\par
El equipo en pleno realizó una revisión exhaustiva de los reportes de Aseguramiento de Calidad (QA). Se analizaron a detalle los casos de prueba (\textit{Test Cases}) ejecutados y los reportes de errores (\textit{Bug Reports}) abarcando todo el ciclo de vida desde el Sprint 1 y 2, hasta el actual Sprint 3. Mediante esta inspección colaborativa, se verificó que las incidencias críticas fueran debidamente trazadas y corregidas, garantizando que el incremento cumple con los criterios de aceptación y los estándares de calidad exigidos por la organización.

\vspace*{0.5cm}
\noindent \textbf{Entrega de Documentación Técnica Integral}\par
Se oficializó la entrega de la documentación técnica completa correspondiente al desarrollo del sistema. El equipo constató la recepción de los manuales y especificaciones arquitectónicas de las capas de \textit{Backend} y \textit{Frontend}, así como el modelado y diccionario de la \textit{Base de Datos}. Adicionalmente, se entregó y validó la documentación pertinente a los flujos de validaciones del sistema y la gestión del control de versiones (\textit{Versioning}). Todo el equipo acordó que estos documentos reflejan fielmente el estado actual del producto, asegurando su futura mantenibilidad y escalabilidad.

\vspace*{0.5cm}
\noindent \textbf{Definición de Fechas de Entrega y Despliegue (Sprint 3)}\par
Con el objetivo de coordinar el cierre del ciclo de manera exitosa, el equipo sometió a consenso el cronograma de finalización del Sprint 3. Tras evaluar la capacidad restante y el estado de los entregables, se acordó por unanimidad fijar la \textbf{fecha de entrega técnica de todos los trabajos del sprint para el día sábado, 9 de mayo, a horas 12:00 (mediodía)}. Así mismo, la ejecución del despliegue oficial (\textit{Deployment}) a los entornos correspondientes quedó programada indudablemente para ese mismo día a las \textbf{18:00 horas}. Todos los miembros expresaron su compromiso absoluto con estos hitos críticos.

\vspace*{0.5cm}
\noindent \textbf{Aprobación Oficial}\par
Al corroborar la recepción conforme de todos los reportes de QA y de la documentación técnica en todas sus áreas, y habiendo consolidado un cronograma preciso y avalado por todos para la finalización y despliegue del tercer sprint, el equipo manifiesta su total conformidad. Se da por aprobada la sesión y se procede con la firma de los presentes.

\newpage
\noindent \textbf{Firmas y Aprobación de los Presentes}
\begin{center}
    \noindent
    \setlength{\tabcolsep}{0pt}
    \begin{tabular}{ccc}
        \espacioFirma{assets/img/firmas/mauricio.jpg}{Mauricio Jaimes Arnez}{13751221} & 
        \espacioFirma{assets/img/firmas/camila.jpg}{Camila Belen Quispe Flores}{13097244} &
        \espacioFirma{assets/img/firmas/jim.jpg}{Jim Gabriel Ariñez Bautista}{9517642} \\[1.0cm]
        \espacioFirma{assets/img/firmas/mateo.jpg}{Mateo Medina Rodríguez}{9453809} &
        \espacioFirma{assets/img/firmas/tomas.jpg}{Tomas Molina Maldonado}{8051701} &
        \espacioFirma{assets/img/firmas/angheline.jpg}{Angheline Gutierrez Guzmán}{13751815} \\[1.0cm]
    \end{tabular}
    \vspace{0.2cm} 
    \begin{tabular}{cc}
        \espacioFirma{assets/img/firmas/diego.jpg}{Juan Diego Pardo Pozo}{12747783}
    \end{tabular}
\end{center}

\vspace{0.5cm}
\noindent \textbf{Fecha de última revisión:} Jueves, 7 de Mayo de 2026